\section{Bilinear maps}
\label{section-bilinear}


\noindent
Let $(X, \mathcal{O}_X)$ be a ringed space.
Let $\mathcal{F}$, $\mathcal{G}$, and $\mathcal{H}$ be $\mathcal{O}_X$-modules.
A {\it bilinear map} $f : \mathcal{F} \times \mathcal{G} \to \mathcal{H}$
of sheaves of $\mathcal{O}_X$-modules is a map of sheaves of sets as indicated
such that for every open $U \subset X$ the induced map
$$
\mathcal{F}(U) \times \mathcal{G}(U) \to \mathcal{H}(U)
$$
is an $\mathcal{O}_X(U)$-bilinear map of modules. Equivalently you can ask
certain diagrams of maps of sheaves of sets commute, imitating the usual
axioms for bilinear maps of modules. For example, the axiom
$f(x + y, z) = f(x, z) + f(y, z)$ is represented by the commutativity
of the diagram
$$
\xymatrix{
\mathcal{F} \times \mathcal{F} \times \mathcal{G}
\ar[rrr]_{(f \circ \text{pr}_{13}, f \circ \text{pr}_{23})}
\ar[d]_{(+ \circ \text{pr}_{12}, \text{pr}_3)} & & &
\mathcal{H} \times \mathcal{H}
\ar[d]^{+} \\
\mathcal{F} \times \mathcal{G} \ar[rrr]^f & & &
\mathcal{H}
}
$$
Another characterization is this: if
$f : \mathcal{F} \times \mathcal{G} \to \mathcal{H}$
is a map of sheaves of sets and
it induces a bilinar map of modules on stalks for all points of $X$, then
$f$ is a bilinear map of sheaves of modules. This is true as you can test
whether local sections are equal by checking on stalks.

\medskip\noindent
Let $\text{Mor}( - , - )$ denote morphisms in the category of sheaves of
sets on $X$. Another characterization of a bilinear map is this:
a map of sheaves of sets
$f : \mathcal{F} \times \mathcal{G} \to \mathcal{H}$
is bilinear if given any sheaf of sets $\mathcal{S}$ the rule
$$
\text{Mor}(\mathcal{S}, \mathcal{F}) \times
\text{Mor}(\mathcal{S}, \mathcal{G}) \to
\text{Mor}(\mathcal{S}, \mathcal{H}),\quad
(a, b) \mapsto f \circ (a \times b)
$$
is a bilinear map of modules over the ring
$\text{Mor}(\mathcal{S}, \mathcal{O}_X)$.
We don't usually take this point of view as it is easier to think about
sets of local sections and it is clearly equivalent.

\medskip\noindent
Finally, here is yet another way to say the definition:
$\mathcal{O}_X$ is a ring object in the category of sheaves of sets
and $\mathcal{F}$, $\mathcal{G}$, $\mathcal{H}$ are module objects
over this ring. Then a bilinear map can be defined for module objects
over a ring object in any category.
To formulate what is a ring object and what is a module object
over a ring object, and what is a bilinear map of such in a
category it is pleasant (but not strictly necessary)
to assume the category has finite products; and
this is true for the category of sheaves of sets.







